% TOP_PROBLEM: {"id": 2980, "title": "Kirby Problem 4.104", "category_id": 7, "category": {"id": 7, "name": "topology", "display_name": "Topology", "description": "Properties preserved under continuous deformations.", "slug": "topology", "order_index": 7, "created_at": "2026-07-31T15:26:25.671Z"}, "status": "open"}
% FIRST_POSTED: null
% ENTRY_KIND: statement_only
% Imported from: batch13.tex; UnsolvedMath ID 2980
\documentclass[11pt]{article}
\usepackage[margin=25mm]{geometry}
\usepackage{fontspec}
\setmainfont{Latin Modern Roman}
\usepackage{amsmath,amssymb,mathtools}
\usepackage{unicode-math}
\setmathfont{Latin Modern Math}
\usepackage{xurl,hyperref}
\hypersetup{colorlinks=true,urlcolor=blue}
\setcounter{secnumdepth}{0}
\setlength{\parindent}{0pt}
\setlength{\parskip}{5pt}
\setlength{\emergencystretch}{3em}
\newcommand{\sref}[2]{\hyperlink{source:#1:#2}{[S#2]}}
\newcommand{\cataloguescope}[1]{\paragraph{Recorded target.}#1}
\begin{document}
% UnsolvedMath ID 2980; source catalogue code KP-4.104
\setcounter{section}{2979}
\section{Kirby Problem 4.104}\label{2980}
{\small\sffamily UnsolvedMath ID 2980 · Topology}
\subsection{Definitions and mathematical statement}

\paragraph{Shared notation.}$\mathbb N=\{1,2,\ldots\}$, and $\mathbb Z,\mathbb Q,\mathbb R,\mathbb C$ denote the integers, rationals, reals and complex numbers. Any other conventions are specified locally.


Let $B^4\subset\mathbb C^2$ be the unit ball with $\omega_{\mathrm{std}}=\sum_{j=1}^2dx_j\wedge dy_j$. On $S^3=\partial B^4$ use the standard cooriented contact structure $\xi_{\mathrm{std}}=\ker\alpha_{\mathrm{std}}$, where
\[
\alpha_{\mathrm{std}}=\tfrac12\sum_{j=1}^2(x_j\,dy_j-y_j\,dx_j)\big|_{S^3}.
\]
A positive transverse link is an oriented smooth link whose tangent vector has positive $\alpha_{\mathrm{std}}$ value. Use a fixed transverse link type $L$, meaning equivalence by an isotopy through positive transverse links.

A complex curve in the ball is a nonsingular properly embedded complex one-dimensional analytic surface, smooth up to its transverse boundary on $S^3$. A smooth isotopy is a family of proper smooth embedded surfaces; a complex isotopy requires every member to be a complex curve; a symplectic isotopy requires $\omega_{\mathrm{std}}$ to restrict to a positive area form on each member. The boundaries of the compared complex curves represent the same transverse link type $L$, and need not be one pointwise fixed geometric analytic boundary curve. The original examples cited by K3 explicitly use transversely isotopic endpoint boundaries, as explained in Source S3. Smooth isotopy here is ordinary isotopy through proper smooth embeddings; no requirement that every intermediate boundary remain transverse is added. \sref{2980}{3}
\paragraph{Statement recorded in the supplied source.}
Does there exist a transverse link type $L$ with two complex curves in $B^4$ that are smoothly isotopic but not complexly isotopic? Does there exist such a pair that is smoothly isotopic but not symplectically isotopic? For each version, can there be an infinite family of complex curves with boundary type $L$, all in one smooth isotopy class and pairwise in different complex, respectively symplectic, isotopy classes? The questions concern separation within a smooth isotopy class, rather than curves already distinguished by smooth topology. \sref{2980}{2}
\subsection{Short English statement}
Can complex curves with the same transverse boundary link type be smoothly isotopic while not complexly or symplectically isotopic, and can infinitely many inequivalent curves occur in one smooth isotopy class?
\subsection{Sources}
\begin{description}
\item[\hypertarget{source:2980:1}{S1}] ULAM.AI and the UnsolvedMath Contributors, UnsolvedMath: A Curated Collection of Open Mathematics Problems, record 2980 (KP-4.104). CC BY 4.0. \url{https://huggingface.co/datasets/ulamai/UnsolvedMath}.
\item[\hypertarget{source:2980:2}{S2}] R. İnanç Baykur, Robion C. Kirby and Daniel Ruberman, K3—A New Problem List in Low-Dimensional Topology, Mathematical Surveys and Monographs 295 (2026), author version, Problem 4.104 and Remarks, PDF page 277. \url{https://math.berkeley.edu/sites/default/files/surv-295-ruberman-watermarked-author-pdf.pdf}.
\item[\hypertarget{source:2980:3}{S3}] Kyle Hayden, Exotically knotted disks and complex curves, arXiv:2003.13681v2 (2021), Sections 4.1--4.2, pages 23--24, proof of Theorem A. \url{https://arxiv.org/pdf/2003.13681v2}.
\end{description}
\label{end:2980}
\end{document}
